\documentclass[10pt]{article}
\usepackage[T1]{fontenc}
\usepackage[utf8]{inputenc}
\usepackage[french]{babel}
\usepackage{amsmath,amssymb,mathrsfs,mathpazo}
\usepackage{geometry}
\geometry{a4paper,margin=2cm}
\usepackage{graphicx,array,tabularx,tabularray,colortbl,multirow,multicol,booktabs,enumitem,pifont,ulem,fancybox,diagbox,cancel}
\usepackage{pstricks,pstricks-add,pst-plot,pst-tree,pst-node,pst-text,pst-func,pst-math,pst-bspline,pst-3dplot,pst-eucl,pst-solides3d}
\let\pscalculate\undefined
\usepackage{pst-bezier}
\usepackage{tikz,pgfplots,tkz-tab}
\usetikzlibrary{arrows,arrows.meta,patterns,calc,positioning,trees,decorations.markings,shapes.geometric,shapes.misc,shapes.arrows,intersections,decorations.pathmorphing,decorations.pathreplacing,backgrounds,shadows,babel}
\pgfplotsset{compat=1.18}
\usepgfplotslibrary{fillbetween}
\usepackage[np]{numprint}
\usepackage{listings,makecell,pas-tableur,twoopt}
\let\barre\undefined
\usepackage[pstricks]{tablvar}
\definecolor{darkgreen}{RGB}{0,100,0}
\providecommand{\red}{\color{red}}
\providecommand{\blue}{\color{blue}}
\providecommand{\green}{\color{green}}
\providecommand{\black}{\color{black}}
\providecommand{\white}{\color{white}}
\providecommand{\gray}{\color{gray}}
\providecommand{\cyan}{\color{cyan}}
\providecommand{\magenta}{\color{magenta}}
\providecommand{\yellow}{\color{yellow}}
\providecommand{\violet}{\color{violet}}
\providecommand{\orange}{\color{orange}}
\providecommand{\euro}{\texteuro}
\providecommand{\EUR}{\texteuro}
\providecommand{\ieme}{\textsuperscript{e}}
\providecommand{\degres}{\ensuremath{^\circ}}
\providecommand{\degree}{\ensuremath{^\circ}}
\providecommand{\vv}[1]{\overrightarrow{#1}}
\providecommand{\Coord}[1]{#1}
\providecommand{\decosix}{}
\providecommand{\starredbullet}{\ensuremath{\bullet}}
\pagestyle{empty}
\setlength{\parindent}{0pt}
\usepackage{adjustbox}\DeclareUnicodeCharacter{2236}{:}\DeclareUnicodeCharacter{2190}{\ensuremath{\leftarrow}}
\providecommand{\R}{\mathbb{R}}
\providecommand{\N}{\mathbb{N}}
\providecommand{\Z}{\mathbb{Z}}
\providecommand{\Q}{\mathbb{Q}}
\providecommand{\D}{\mathbb{D}}
\providecommand{\C}{\mathbb{C}}
\providecommand{\e}{{\mathrm{e}}}
\providecommand{\vect}[1]{\overrightarrow{#1}}
\providecommand{\vectt}[1]{\overrightarrow{\mathrm{#1}}}
\providecommand{\vvt}[1]{\overrightarrow{\mathrm{#1}}}
\providecommand{\np}[2][]{{#2~#1}}
\providecommand{\vv}[1]{{\overrightarrow{#1}}}
\renewcommand{\d}{{\mathrm{d}}}
\providecommand{\degres}{{^\circ}}
\providecommand{\degre}{{^\circ}}
\providecommand{\blue}{\color{blue}}
\providecommand{\red}{\color{red}}
\providecommand{\enskip}{\quad}
\providecommand{\llbracket}{\lbrack\!\lbrack}
\providecommand{\rrbracket}{\rbrack\!\rbrack}
\providecommand{\euro}{{\text{€}}}
\providecommand{\no}{{\text{n°}}}
\providecommand{\interval}[2]{{[#1;#2]}}
\providecommand{\intervFO}[2]{{[#1;#2[}}
\providecommand{\dotfill}{\ldots}

\providecommand{\ds}{\displaystyle}
\providecommand{\textding}[1]{#1}
\providecommand{\Oij}{\(\mathrm{O};\vec{\imath},\vec{\jmath}\)}
\providecommand{\Oijk}{\(\mathrm{O};\vec{\imath},\vec{\jmath},\vec{k}\)}
\providecommand{\Ouv}{\(\mathrm{O};\vec{u},\vec{v}\)}
\providecommand{\intervFF}[2]{\ensuremath{[#1;#2]}}
\renewcommand{\euro}{\ifmmode\text{\texteuro}\else\texteuro\fi}
\DeclareUnicodeCharacter{00B0}{\ifmmode^{\circ}\else\textdegree\fi}
\DeclareUnicodeCharacter{20AC}{\ifmmode\text{\texteuro}\else\texteuro\fi}
% Énoncés : Éducation nationale. Transcription et corrigés : APMEP, Denis Vergès et contributeurs. Découpage et classement : CoopMaths, Sébastien Lozano, Éric Elter et contributeurs. Adaptation de la présentation : Maths Explorer.
% CC BY-SA : https://coopmaths.fr/www/about/
% https://coopmaths.fr/alea/static/eam/2026/tex/eam_specialite_2026_06_antilles_2.tex
\begin{document}
\noindent\textit{APMEP / CoopMaths. CC BY-SA. Présentation adaptée par Maths Explorer.}\par\medskip
On s'intéresse à l'ensemble des spectateurs d'un cinéma.

On désigne par le terme \emph{JEUNE} un spectateur âgé de 18 ans ou moins.

On désigne par le terme \emph{ABONNÉ} un spectateur ayant souscrit un abonnement à ce cinéma.

Une étude a permis d'établir les résultats suivants :

\begin{itemize}[label*=\textbullet~]
	\item 60\,\% des spectateurs sont abonnés ;
	\item parmi les spectateurs abonnés, 40\,\% sont jeunes ;
	\item parmi les spectateurs n'étant pas abonnés, 20\,\% sont jeunes.
\end{itemize}

On interroge au hasard un spectateur et on considère les événements suivants :
\begin{itemize}[label*=\textbullet~]
	\item $A$ : \og \emph{le spectateur est abonné} \fg{} ;
	\item $J$ : \og \emph{le spectateur est jeune} \fg.
\end{itemize}

\begin{enumerate}
	\item Reproduire et compléter l'arbre pondéré ci-dessous décrivant la situation.

	
	\begin{center}
		\begin{tikzpicture}[yscale=-1,
			grow = right,
			level distance=2.5cm,
			level 1/.style={sibling distance=1.6cm},
			level 2/.style={sibling distance=0.8cm},
			edge from parent/.style={draw},
			edge from parent path=
			{(\tikzparentnode.east)--(\tikzchildnode.west)},
			every node/.style={fill=white}]
		\node {}
		child {
				node {$A$}
					child { node {$J$} edge from parent node {\dots} }
					child { node {$\overline{J}$} edge from parent node {\dots} }
				edge from parent node {\dots}
				}
		child {
				node {$\overline{A}$}
					child { node {$J$} edge from parent node {\dots} }
					child { node {$\overline{J}$} edge from parent node {\dots} }
				edge from parent node {\dots}
				};
		\end{tikzpicture}













	\end{center}

	\item Calculer la probabilité $P(A \cap J)$ et interpréter le résultat dans le contexte de l'exercice.

	\item Démontrer que $P(J) = 0,32$.

	\item On considère l'affirmation suivante :

	\begin{quote}
		« \textit{Si on interroge un spectateur jeune au hasard, il y a plus d'une chance sur deux qu'il soit abonné.} »
	\end{quote}

	Cette affirmation est-elle exacte ? Justifier.

	\item Le prix d'une place de cinéma est fixé selon les règles suivantes :
	\begin{itemize}[label*=\textbullet~]
		\item un spectateur non abonné paie sa place 5 euros, quel que soit son âge ;
		\item un spectateur abonné et n'étant pas jeune paie sa place 2 euros ;
		\item un spectateur abonné et jeune ne paie pas sa place : c'est gratuit.
	\end{itemize}

	On note $X$ la variable aléatoire qui, à un spectateur choisi au hasard, donne le prix qu'il doit payer pour sa place de cinéma.

	\begin{enumerate}
		\item Déterminer la loi de probabilité de la variable aléatoire $X$.
		\item Calculer l'espérance de la variable aléatoire $X$ et interpréter le résultat dans le contexte de l'exercice.
	\end{enumerate}
\end{enumerate}

\end{document}
